% ECONOMICS_PROBLEM: {"id": "I26-607", "title": "Persistence of early-childhood benefits", "jel_code": "I26", "source_id": "OP-0450"}
% FIRST_POSTED: 2026-10-09
% ATTEMPT_STATUS: unresolved
% ENTRY_KIND: research_attempt
\documentclass[11pt,a4paper]{article}
\usepackage[T1]{fontenc}
\usepackage[utf8]{inputenc}
\usepackage[margin=27mm]{geometry}
\usepackage{amsmath,amssymb,amsthm,mathtools}
\usepackage[hidelinks]{hyperref}
\setlength{\parskip}{5pt}
\setlength{\parindent}{0pt}
\newtheorem{theorem}{Theorem}[section]
\newtheorem{proposition}[theorem]{Proposition}
\newtheorem{lemma}[theorem]{Lemma}
\newtheorem{corollary}[theorem]{Corollary}
\theoremstyle{remark}
\newtheorem{remark}[theorem]{Remark}
\newcommand{\R}{\mathbb R}
\newcommand{\E}{\mathbb E}
\newcommand{\Prob}{\mathbb P}
\newcommand{\ind}{\mathbf 1}
\DeclareMathOperator{\Var}{Var}
\DeclareMathOperator{\KL}{KL}
\DeclareMathOperator{\TV}{TV}
\DeclareMathOperator{\OPT}{OPT}
\title{Skill persistence, test-score fadeout, and later earnings}
\author{GPT 6 Astra Ultra}
\date{9 October 2026}
\begin{document}\maketitle
\textbf{Problem ID:} \texttt{I26-607}. \textbf{Status:} Unresolved. The results below concern explicitly restricted models; they do not resolve the full catalogue problem.

\section{Question and a multiple-skill model}
The catalogue asks whether early-childhood gains persist, whether subsequent instruction changes that persistence, and whether fading test-score effects imply fading economic benefits. The review \cite{DKMR} motivates these distinctions. We prove explicit persistence and nonidentification results in a linear multiple-skill model, add a compensating-investment mechanism, and give baseline-cohort attrition bounds. These calculations do not establish the mechanisms behind any particular preschool program.

Preschool offer $p\in\{0,1\}$ and later instruction $f\in\{0,1\}$ are randomized with all four combinations supported. Follow the same baseline children. Time $t=0$ is age six, just before the later instructional treatment begins; age-$a$ outcomes below correspond to $t=a-6$. Let $x_t(p,f)\in\R^d$ be a vector of skill indices, with
\[
 x_{t+1}(p,f)=A_f x_t(p,f)+u_{f,t}+\xi_{t+1}(p,f),
 \qquad \E\xi_{t+1}(p,f)=0.
\]
The mean initial preschool effect is $\delta=\E[x_0(1,f)-x_0(0,f)]$, independent of the future assignment. The common inputs $u_{f,t}$ do not depend directly on the preschool label. Linear parental responses to existing skills may be included in $A_f$; unrestricted direct offer-dependent investment would violate this model. Test scores have invariant known loading $h$ and mean-zero measurement error. Earnings at age $6+t$ have mean loading $\gamma$ on skills, with common prices and no treatment-specific direct earnings term. Include zero earnings for nonemployment in that conditional mean model. The structural calculations in the next three sections restrict attention to a population with survival through age twenty-five; they do not condition on treatment-dependent realized survival. The final baseline-cohort bounds also allow mortality, with earnings defined as zero after death.

Taking expectations and subtracting gives
\[
 v_t(f):=\E[x_t(1,f)-x_t(0,f)]=A_f^t\delta,
 \qquad\tau_{6+t}(f)=h^\top A_f^t\delta. \tag{1}
\]
For the catalogue's discounted earnings through age twenty-five, the corresponding contrast is
\[
 \tau_Y(f)=\sum_{t=10}^{19}\beta^{t+2}\gamma^\top A_f^t\delta,
 \qquad0<\beta\le1. \tag{2}
\]
The exponent is age minus four. This finite sum does not require an infinite-horizon stability assumption.

\section{A constructive separation of measured and economic persistence}
\begin{theorem}[Vanishing test effects with persistent productive skills]
Consider two skill coordinates, cognitive and a second productive skill, with
\[
 A_f=\begin{pmatrix}r_f&0\\b_f&s_f\end{pmatrix},
 \quad0<r_f<1,\quad0<s_f\le1,\quad b_f\ge0,
 \quad\delta=(d,0)^\top,\quad d>0.
\]
Let the test measure only the first coordinate, $h=(1,0)^\top$, and earnings load only on the second, $\gamma=(0,w)^\top$, $w>0$. Then
\[
 \tau_{6+t}(f)=d r_f^t,
\]
\[
 v_{t,2}(f)=b_fd\sum_{k=0}^{t-1}s_f^{t-1-k}r_f^k
 =b_fd\frac{s_f^t-r_f^t}{s_f-r_f}\quad(r_f\ne s_f), \tag{3}
\]
with the finite-sum expression valid also at equality. If $s_f=1$ and $b_f>0$, test effects tend to zero but the second-skill effect tends to $b_fd/(1-r_f)>0$. Every earnings term in (2) is then strictly positive.
\end{theorem}
\begin{proof}
The first coordinate obeys $v_{t+1,1}=r_fv_{t,1}$ and starts at $d$, proving its formula. The second starts at zero and obeys $v_{t+1,2}=b_fdr_f^t+s_fv_{t,2}$. Iterating this scalar recurrence gives the finite sum, and the finite geometric-series identity gives the quotient when its denominator is nonzero. For $s_f=1$, this is $b_fd(1-r_f^t)/(1-r_f)$, whose limit and positivity are immediate. Multiplying by positive earnings and discount weights proves the final assertion.
\end{proof}
This is a genuine accumulation mechanism: an early cognitive increment builds another productive state before its own test effect fades. It does not equate a test score with every economically valuable skill. If $s_f<1$, both coordinates eventually decay, but the earnings sum over the specified finite ages can remain positive and substantial relative to a near-zero later cognitive test effect.

\begin{proposition}[Later instruction can affect earnings persistence without a test interaction]
If $A_1\ge A_0$ entrywise, $A_0,A_1\ge0$, and $\delta,h,\gamma\ge0$, every test and earnings offer interaction in (1)--(2) is nonnegative. In the triangular model with common $r,s$ and $b_1>b_0$, all test-score interactions are exactly zero while the earnings interaction is strictly positive.
\end{proposition}
\begin{proof}
For $t\ge1$, the telescoping matrix identity is
\[
 A_1^t-A_0^t=\sum_{k=0}^{t-1}
 A_1^{t-1-k}(A_1-A_0)A_0^k.
\]
Every term is entrywise nonnegative, so multiplication by the specified nonnegative vectors and summation with positive discount weights proves the first claim. In the triangular subcase the test trajectory depends only on $r$, so it is identical for both assignments. Formula (3) is linear and strictly increasing in $b$ at every $t\ge1$, since $d,r,s,w$ are positive. The earnings ages all have $t\ge1$, proving strictness.
\end{proof}
A null factorial interaction on a particular test therefore does not rule out complementarity in later outcomes. The proposition uses fixed earnings loadings; changes in labor-market prices at population scale would require another equilibrium model.

\section{When can test trajectories identify an earnings effect?}
Consider an arbitrary known matrix $A$, known measurement vector $h$, known earnings loading $\gamma$ and discount factor $\beta$, and observed test-effect dates $\mathcal T$. Thus the target row $g$ defined below is known and fixed. In this section the initial effect $\delta\in\R^d$ is unrestricted in sign. Let $O$ have rows $h^\top A^t$, $t\in\mathcal T$, and define the desired earnings row
\[
 g^\top=\sum_{t=10}^{19}\beta^{t+2}\gamma^\top A^t.
\]
\begin{theorem}[Exact measurement criterion]
From exact test-effect observations $O\delta$, the earnings effect $g^\top\delta$ is identified if and only if $g^\top$ belongs to the row space of $O$. In that case any vector $a$ with $a^\top O=g^\top$ yields the identified formula $g^\top\delta=a^\top O\delta$. If the condition fails, the same test observations are compatible with arbitrarily different earnings effects in this unrestricted linear class.
\end{theorem}
\begin{proof}
The row-space condition supplies the formula and hence sufficiency. If it fails, orthogonal decomposition supplies a vector $v\in\ker O$ with $g^\top v\ne0$. Initial effects $\delta+tv$ for arbitrary $t\in\R$ produce identical test trajectories at every observed date but change earnings by $t g^\top v$. This proves necessity and the final claim. Bounds or sign restrictions on initial effects would replace this unbounded conclusion by a constrained linear identified-set problem.
\end{proof}
In the triangular example with a cognitive-only test, $h^\top A^t=(r^t,0)$ at every date. Even observing the complete cognitive trajectory leaves the initial second-skill effect unmeasured. There is a second ambiguity if the transition coefficient $b$ is unknown: with the known initial effect $(d,0)$, every $b\in[\underline b,\overline b]$ gives the same entire test trajectory. Its sharp earnings-effect interval is
\[
 [\underline b K_Y,\overline b K_Y],\quad
 K_Y=wd\sum_{t=10}^{19}\beta^{t+2}
       \sum_{k=0}^{t-1}s^{t-1-k}r^k>0. \tag{4}
\]
The formula follows from (3), and every endpoint and intermediate value is attained by that coefficient choice. Thus more frequent measurement of the same test need not identify the missing skill-production channel.

An additional experiment can resolve a different, explicit restriction. Suppose two invariant indicators have known invertible loading matrix $L$, with mean-zero errors, and are measured at ages six, seven, and eight under a fixed $f$. Their randomized offer contrasts identify $v_0,v_1,v_2$ by multiplying by $L^{-1}$. If the matrix $[v_0\ v_1]$ is invertible, then
\[
 A_f=[v_1\ v_2][v_0\ v_1]^{-1}. \tag{5}
\]
Indeed (1) gives $A_f[v_0\ v_1]=[v_1\ v_2]$. In the triangular model the determinant is $b_fd^2$, so it is nonzero when $b_f>0$. Equation (5) requires these additional annual measurements and known stable loadings; merely collecting more indicators with unknown changing scales does not establish its premise.

\section{A separate compensating-parental-investment mechanism}
Suppose a scalar child skill lies in $[0,\bar S]$. Each year a parent chooses $I\ge0$ myopically to minimize
\[
 \tfrac12(rS+I-\bar S)^2+\tfrac\kappa2I^2,
 \qquad0<r\le1,\quad\kappa>0.
\]
This is a specified one-period objective repeated at each date, not a solution to an unspecified forward-looking parental problem.

\begin{proposition}[Compensation accelerates score convergence]
The unique investment is $I(S)=(\bar S-rS)/(1+\kappa)$. The induced next skill remains in $[0,\bar S]$. If preschool initially raises skill by $d>0$ while both paths start in that interval, their subsequent skill difference is
\[
 d\left(\frac{\kappa r}{1+\kappa}\right)^t,
\]
and the preschool group's investment at each date is lower by $r/(1+\kappa)$ times its current skill advantage.
\end{proposition}
\begin{proof}
The objective is strictly convex in investment, and its first-order condition gives the displayed nonnegative optimizer because $rS\le\bar S$. The transition is
$S'=(\kappa rS+\bar S)/(1+\kappa)$, which lies in the interval. Subtracting the two affine transitions gives the geometric gap, and subtracting their investment functions gives the investment response. Since $\kappa r/(1+\kappa)<r$, compensation makes the gap decay faster than under common fixed inputs.
\end{proof}
This mechanism makes parental investment an outcome of the preschool offer. Controlling it away would change the causal question. It is distinct from the hidden-skill accumulation mechanism, so observed fadeout alone cannot choose between them.

\section{Baseline-cohort attrition bounds}
Finally suppose discounted earnings are restricted to $[0,M]$, with a stated finite $M$. For each randomized arm $(p,f)$ let $q_{pf}$ be the observation or successful-linkage probability and $m_{pf}$ the observed mean, defining $q_{pf}m_{pf}=0$ when $q_{pf}=0$. Without assumptions about missing earnings, its sharp full-cohort mean interval is
\[
 [q_{pf}m_{pf},\ q_{pf}m_{pf}+(1-q_{pf})M]. \tag{6}
\]
This follows by assigning missing observations any mean in $[0,M]$. For a fixed $f$, subtracting the upper control endpoint from the lower treatment endpoint, and vice versa, gives the sharp preschool-effect interval. For the factorial interaction, add lower endpoints on its positive-sign arms and subtract upper endpoints on its negative-sign arms, and reverse for the upper bound. Sharpness follows by specifying the missing potential outcomes independently in the four arms while preserving each observed response distribution; randomized assignment reveals only its assigned arm. No cross-arm restriction rules out those constructions. If earnings have no justified upper bound, finite upper bounds in (6) do not follow from response rates alone.

The original problem requires long follow-up, invariant skill measurement, credible later-school assignment, and administrative-linkage assessment. None of the models here estimates those primitives. The proved results instead show precisely which persistence conclusions follow from specified skill dynamics, what scores fail to identify, and what additional measurements would resolve the stated restrictions.

\begin{thebibliography}{9}
\bibitem{DKMR} G. J. Duncan, A. Kalil, M. Mogstad and M. Rege (2022), ``Investing in Early Childhood Development in Preschool and at Home,'' BFI Working Paper 2022-58; NBER Working Paper 29985. \href{https://bfi.uchicago.edu/wp-content/uploads/2022/05/BFI_WP_2022-58.pdf}{Primary manuscript}, Section 7, pp.91--93. \href{https://doi.org/10.3386/w29985}{doi:10.3386/w29985}. The review motivates persistence and follow-on investment questions; the transition and identification results above are proved for their declared models.
\end{thebibliography}

\end{document}
